\chapter{Pipeline phân tích nhật ký}

\section{Hợp đồng đầu vào và xử lý lỗi}

Parser nhận CSV hoặc JSON và chuẩn hóa về Event có các trường
\texttt{event\_id}, timestamp, event type, user, IP, resource, action, status,
request, details và \texttt{expected\_label}. Ba event type hợp lệ là
authentication, access và authorization. Header thiếu, loại file không hỗ trợ
hoặc schema cấp file sai bị từ chối; các dòng có timestamp sai được gom thành
\texttt{RowError} để dòng hợp lệ còn lại vẫn tiếp tục được xử lý.

Trường \texttt{expected\_label} chỉ có trong dataset đánh giá. Detection engine
không đọc nó khi tạo dự đoán, nhờ đó việc tính metric về sau không làm rò nhãn
vào quá trình phát hiện.

\section{Rule nền và chấm điểm}

\texttt{SQLInjectionRule} dùng các regex được compile sẵn và dừng tại pattern
đầu tiên khớp với mỗi event. \texttt{PasswordGuessingRule} nhóm authentication
failure theo cặp user--IP, sắp xếp thời gian và tạo finding khi số lần thất bại
đạt ngưỡng trong cửa sổ cấu hình. \texttt{UnauthorizedAccessRule} chỉ xét
access/authorization có user, gọi RBAC repository và kết hợp quyền thực tế với
status denied trong log.

Risk scorer đặt điểm nền 45 cho SQL Injection, 40 cho password guessing và 50
cho unauthorized access. Điểm cộng bị giới hạn cho confidence, số lần lặp và
trạng thái RBAC từ chối; tổng luôn nằm trong khoảng 0--100. Severity là Medium
từ 30, High từ 60 và Critical từ 85 theo \texttt{config/default.toml}.

Với finding nền, công thức được hiện thực là:
\[
S=\min\left(100,\;B+\operatorname{round}(20c)
  +\min(15,3\max(0,n-1))+15I_{denied}\right),
\]
trong đó $B$ là điểm nền, $c\in[0,1]$ là confidence do rule gán, $n$ là số
Event lặp lại và $I_{denied}$ bằng 1 khi repository xác nhận RBAC từ chối.
Ví dụ, một unauthorized access với confidence 0,9 và quyết định denied nhận
$50+18+15=83$ điểm, tương ứng High. Công thức giúp giải thích thứ tự ưu tiên,
nhưng không được hiệu chỉnh từ tần suất thực tế; 83 điểm không có nghĩa xác
suất tấn công là 83\%.

\section{Context-risk và incident}

Khi bật \texttt{--context-risk}, \texttt{BehaviorProfiler} duyệt Event theo
thời gian để nhận diện ngoài giờ, IP mới của người dùng, tài nguyên chưa từng
truy cập thành công và ba lần từ chối trong cửa sổ cấu hình. Đây là heuristic
phục vụ triage, không phải một bộ phân loại được đánh giá độc lập.

Finding ngữ cảnh dùng $S_c=\min(100,25+\operatorname{round}(20c)+b)$, với
$b$ là bonus theo loại tín hiệu. Cách tách công thức tránh làm context signal
trông như một finding nền đã được rule xác nhận. Các ngưỡng severity và bonus
là giả định thiết kế công bố trong cấu hình; muốn dùng trong vận hành phải hiệu
chỉnh trên workload và chi phí triage của tổ chức.

Sau scoring, các alert cùng user--IP được phân vùng theo khoảng thời gian. Mỗi
incident giữ alert ID, event ID, loại rủi ro, tín hiệu ngữ cảnh và evidence có
thứ tự ổn định. Điểm incident là điểm alert lớn nhất cộng session-chain bonus,
sau đó bị giới hạn 100. Luồng này được minh họa ở Hình~\ref{fig:event-alert}.

\begin{figure}[htbp]
\centering
\begin{tikzpicture}[
  node distance=9mm and 8mm,
  box/.style={draw, rounded corners, align=center, minimum height=11mm, minimum width=28mm, fill=blue!6},
  optional/.style={draw, rounded corners, align=center, minimum height=11mm, minimum width=28mm, fill=orange!12},
  flow/.style={-Latex, thick}
]
\node[box] (event) {CSV / JSON\\Event};
\node[box, right=of event] (parser) {Parser\\validation};
\node[box, right=of parser] (rules) {Rule engine\\RBAC check};
\node[box, right=of rules] (alert) {Alert\\score + severity};
\node[optional, below=of rules] (context) {Context profiler\\IP, giờ, resource, denial};
\node[optional, below=of alert] (incident) {Incident\\grouping + triage};
\draw[flow] (event) -- (parser);
\draw[flow] (parser) -- (rules);
\draw[flow] (rules) -- (alert);
\draw[flow] (parser) -- (context);
\draw[flow] (context) -- (alert);
\draw[flow] (alert) -- (incident);
\end{tikzpicture}

\caption{Luồng từ Event sang Alert và Incident trong chế độ context-risk.}
\label{fig:event-alert}
\end{figure}

\section{Artifact đầu ra}

Mỗi lần phân tích sinh \texttt{alerts.csv}, \texttt{alerts.json} và
\texttt{run\_metadata.json}. Khi context-risk được bật, reporter còn tạo file
finding ngữ cảnh và hai file incident ở cả CSV lẫn JSON. Reporter ghi tệp tạm
rồi replace, tránh để lại artifact dở dang khi ghi lỗi. Lệnh evaluate nhận
alerts và Event có nhãn để sinh \texttt{metrics.json}; khi có
\texttt{--manifest}, nó còn tính metric theo campaign. Người viết báo cáo không
chỉnh sửa thủ công các tệp này. Manifest giữ nguồn gốc, phân tầng dữ liệu, hard
negative, campaign ID và commit đóng băng rule để tránh nhầm dataset với
artifact được tạo sau phân tích.
